% =============================================================================
% SC 2026 Artifact Description (AD) Appendix
% Template: SC26 Reproducibility Initiative (github.com/jennfshr/sc26-repro)
% =============================================================================
\documentclass[conference]{IEEEtran}

\usepackage{cite}
\usepackage{amsmath,amssymb,amsfonts}
\usepackage{graphicx}
\usepackage{textcomp}
\usepackage{xcolor}
\usepackage{booktabs}
\usepackage{hyperref}

% --- SC26 Reproducibility macros (inline to avoid tagging.sty dependency) ---
\newcounter{artifactCounter}

\newcommand{\newartifact}{%
    \stepcounter{artifactCounter}%
    \subsection{Computational Artifact $A_{\theartifactCounter}$}%
}

\newcommand{\artrel}{\subsection*{Relation To Contributions}}
\newcommand{\artexp}{\subsection*{Expected Results}}
\newcommand{\arttime}{\subsection*{Expected Reproduction Time (in Minutes)}}
\newcommand{\artin}{\subsection*{Artifact Setup (incl.\ Inputs)}}
\newcommand{\artinpart}[1]{\subsubsection*{#1}}
\newcommand{\artcomp}{\subsection*{Artifact Execution}}
\newcommand{\artout}{\subsection*{Artifact Analysis (incl.\ Outputs)}}

\newcommand{\apppart}[1]{%
\vspace{2ex}
\begin{center}\begin{large}\textbf{#1}\end{large}\end{center}\par
}

\newcommand{\appendixAD}{\apppart{Artifact Description (AD)}}

\hypersetup{
    colorlinks=true,
    linkcolor=black,
    citecolor=black,
    urlcolor=blue,
}

\begin{document}

\twocolumn[%
{\begin{center}
\Huge
Appendix: Artifact Description/Artifact Evaluation
\end{center}}
]

%%%%%%%%%%%%%%%%%%%%%%%%%%%%%%%%%%%%%%%%%%%%%%%%%%%%%
%  AD Appendix
%%%%%%%%%%%%%%%%%%%%%%%%%%%%%%%%%%%%%%%%%%%%%%%%%%%%%

\appendixAD

\section{Overview of Contributions and Artifacts}

\subsection{Paper's Main Contributions}

\begin{description}
\item[$C_1$] \textbf{ML multi-label detection.}
A biquality learning framework that trains on 187 benchmark ground-truth samples
and 91{,}807 heuristic-labeled production logs to detect 8 I/O bottleneck
dimensions, achieving Micro-F1\,=\,0.923 ($\pm$0.004, 5 seeds) on a held-out
436-sample benchmark test set---2.4$\times$ higher than Drishti heuristics.

\item[$C_2$] \textbf{Benchmark Knowledge Base (KB).}
A construction-verified KB of 623 entries from 6 benchmark suites
(IOR, mdtest, DLIO, custom, h5bench, HACC-IO) across 48 scenario
configurations, mapping \{Darshan signature $\leftrightarrow$ code cause
$\leftrightarrow$ validated fix\}.

\item[$C_3$] \textbf{SHAP-structured LLM recommendations.}
An LLM pipeline (IOPrescriber) that uses ML predictions and per-label
SHAP attributions as structured context for RAG-grounded code-level
recommendations, achieving 1.000 groundedness (Claude, Llama) on 12 workloads.

\item[$C_4$] \textbf{Closed-loop validation.}
End-to-end validation where LLM-generated code fixes are executed on HPC
hardware, yielding measured speedups of 7.2$\times$ geometric mean
(range 3.6--39.0$\times$) across 4 IOR workload pairs.

\item[$C_5$] \textbf{1.37M production dataset analysis.}
Feature extraction and preprocessing pipeline for 1.37M Darshan logs from
ALCF Polaris (DOI:~10.5281/zenodo.15052603), the first I/O bottleneck
detection study on a GPU-accelerated leadership computing facility.
\end{description}

\subsection{Computational Artifacts}

\begin{description}
\item[$A_1$] Source code, trained models, and configuration files.\\
Repository: \url{https://github.com/[anonymous]}
\item[$A_2$] Production Darshan dataset (1.37M logs).\\
DOI: \url{https://doi.org/10.5281/zenodo.15052603}
\item[$A_3$] Benchmark ground-truth dataset and Knowledge Base.\\
Included in $A_1$ under \texttt{data/processed/benchmark/} and
\texttt{data/knowledge\_base/}.
\end{description}

\vspace{1ex}
\begin{center}
\begin{tabular}{rll}
\toprule
Artifact ID & Contributions & Related \\
            & Supported     & Paper Elements \\
\midrule
$A_1$ & $C_1$, $C_3$, $C_4$ & Tables~1--4 \\
      &                      & Figures~3--8 \\
\midrule
$A_2$ & $C_5$ & Table~1 \\
      &       & Figures~1--2 \\
\midrule
$A_3$ & $C_2$, $C_4$ & Tables~2--3 \\
      &              & Figures~4, 7--8 \\
\bottomrule
\end{tabular}
\end{center}

%%%%%%%%%%%%%%%%%%%%%%%%%%%%%%%%%%%%%%%%%%%%%%%%%%%%%%%%
\section{Artifact Identification}
%%%%%%%%%%%%%%%%%%%%%%%%%%%%%%%%%%%%%%%%%%%%%%%%%%%%%%%%

% ===========================================================
%  A1: Source code, models, configs
% ===========================================================
\newartifact

\artrel
Artifact $A_1$ contains all source code for the ML training pipeline
($C_1$), SHAP attribution and LLM recommendation pipeline ($C_3$),
closed-loop validation scripts ($C_4$), and trained model weights.
It enables reproduction of every table and figure in the paper.

\artexp
\begin{itemize}
\item ML detection: XGBoost Micro-F1 $\geq$ 0.915 on the 436-sample
benchmark test set (5-seed mean: 0.923 $\pm$ 0.004).
\item Per-label F1 scores matching Table~2 within reported confidence
intervals.
\item LLM groundedness $\geq$ 0.95 for Claude and Llama models.
\item SHAP feature importance rankings consistent with Figure~5.
\end{itemize}

\arttime
\begin{itemize}
\item \textbf{Setup:} 15 min (conda environment creation, data download).
\item \textbf{Execution:} 60 min total. ML training: $\sim$30 min (5 seeds
$\times$ 5 models). LLM inference: $\sim$20 min (12 workloads $\times$ 3
providers, cached). SHAP analysis: $\sim$10 min.
\item \textbf{Analysis:} 10 min (figure and table generation).
\end{itemize}

\artin

\artinpart{Hardware}
Training and inference were performed on NCSA Delta CPU nodes:
AMD EPYC 7763 (Milan), 128 cores per node, 256\,GB DDR4 RAM.
No GPU is required. The minimum configuration for reproduction is
16 CPU cores and 16\,GB RAM.
Benchmarks and closed-loop validation require access to a Lustre
parallel filesystem and SLURM job scheduler.

\artinpart{Software}
\begin{itemize}
\item Python 3.9 (conda environment specified in \texttt{environment.yml})
\item XGBoost 2.1.4, LightGBM 4.5.0, scikit-learn 1.6.1
\item SHAP 0.49.1, Cleanlab 2.7.1, pandas 2.3.3, NumPy 1.26.4
\item PyDarshan 3.5.0 (Python bindings for libdarshan-util)
\item LLM access: OpenRouter API (\texttt{openai} Python package) for
Claude, GPT-4o, Llama~3.1 70B. API key required for live inference;
cached outputs included for offline reproduction.
\item sentence-transformers 2.2+, faiss-cpu 1.7+ (RAG retrieval)
\end{itemize}

\artinpart{Datasets / Inputs}
\begin{itemize}
\item \textbf{Production features:} Pre-extracted in
\texttt{data/processed/production/} (131{,}151 rows $\times$ 195 columns,
140\,MB Parquet). Raw Darshan logs available at DOI:~10.5281/zenodo.15052603.
\item \textbf{Benchmark ground truth:} 623 labeled samples in
\texttt{data/processed/benchmark/} with labels derived from benchmark
configuration parameters (construction-based, not Darshan-derived).
Split: 187 development / 436 test via iterative stratification.
\item \textbf{Knowledge Base:} 623 JSON entries in
\texttt{data/knowledge\_base/} mapping Darshan signatures to code causes
and validated fixes.
\item \textbf{LLM cache:} Pre-computed LLM outputs in
\texttt{data/llm\_cache/} for offline reproduction without API keys.
\end{itemize}

\artinpart{Installation and Deployment}
\begin{verbatim}
git clone <repo-url> && cd SC_2026
conda env create -f environment.yml
conda activate sc2026
python -c "import xgboost, lightgbm, \
  shap, cleanlab; print('OK')"
\end{verbatim}
No compilation step is needed; all dependencies are pure Python or
pre-built wheels. See \texttt{INSTALL.md} for detailed instructions.

\artcomp
The experiment workflow consists of five tasks with the following
dependencies:

$T_1 \rightarrow T_2 \rightarrow T_3 \rightarrow T_4 \rightarrow T_5$

\begin{description}
\item[$T_1$:] \textbf{Feature extraction and preprocessing.}
Parse Darshan logs, extract 186 features (136 raw + 8 indicators +
39 derived + 3 metadata), apply group-specific normalization (log1p +
RobustScaler for heavy-tailed, none for bounded), temporal
train/val/test split (70/15/15).
Script: \texttt{scripts/reproduce\_all.sh --step 1-4}.
\item[$T_2$:] \textbf{ML training (biquality learning).}
Train 5 model families (XGBoost, LightGBM, RF, MLP, Logistic Regression)
with 5 random seeds each. Biquality weight\,=\,100 for ground-truth
samples. Hyperparameters in \texttt{configs/training.yaml}.
Script: \texttt{scripts/reproduce\_all.sh --step 6}.
\item[$T_3$:] \textbf{SHAP attribution.}
Compute SHAP values for the best model (XGBoost) on the test set.
Generate per-label feature importance plots.
Script: \texttt{scripts/reproduce\_all.sh --step 7}.
\item[$T_4$:] \textbf{LLM recommendation (IOPrescriber).}
For each of 12 representative workloads: ML predicts bottleneck labels,
SHAP identifies contributing features, RAG retrieves matching KB entries,
LLM generates code-level fix recommendations.
Script: \texttt{scripts/reproduce\_all.sh --step 8}.
\item[$T_5$:] \textbf{Figure and table generation.}
Produce all paper figures (30 total) and LaTeX tables from result JSON
files.
Script: \texttt{scripts/reproduce\_all.sh --step 9}.
\end{description}

\textbf{Quick reproduction} (skip raw feature extraction, use
pre-processed data):
\begin{verbatim}
bash scripts/reproduce_all.sh --quick
\end{verbatim}
Expected wall-clock time: $\sim$20\,min on 16+ CPU cores.

\textbf{Statistical parameters:}
5 random seeds (42, 123, 456, 789, 1024) for all ML experiments.
Bootstrap 95\% confidence intervals with 10{,}000 resamples on the
test set. LLM inference at temperature\,=\,0 for determinism.

\artout
\begin{itemize}
\item \texttt{results/final\_metrics.json}: Per-model, per-seed, per-label
F1, precision, recall, Hamming loss, and subset accuracy. Reproduces
Tables~1--2.
\item \texttt{results/phase2/shap/}: SHAP summary plots, per-label
beeswarm plots. Reproduces Figures~5--6.
\item \texttt{results/llm\_evaluation/}: LLM groundedness scores,
latency, token counts per provider. Reproduces Table~3.
\item \texttt{paper/figures/}: All 30 paper figures as PDF files.
\item \texttt{paper/tables/}: LaTeX table source files.
\end{itemize}

% ===========================================================
%  A2: Production Darshan dataset
% ===========================================================
\newartifact

\artrel
Artifact $A_2$ is the 1.37M Darshan log dataset from ALCF Polaris
($C_5$). It serves as the production data source from which 131{,}151
cleaned samples and 91{,}807 heuristic-labeled training samples are derived.

\artexp
After parsing and cleaning, the dataset yields 131{,}151 rows with
195 columns (156 raw + 39 derived features). Temporal split produces
91{,}807 training, 19{,}672 validation, and 19{,}672 test samples.
The heuristic label distribution matches the rates reported in
Table~1 of the paper (e.g., access\_granularity: 30.7\%,
healthy: 54.2\%).

\arttime
\begin{itemize}
\item \textbf{Setup:} 30 min (download 1.37M files from Zenodo, $\sim$50\,GB).
\item \textbf{Execution:} 120 min (batch feature extraction on 128 cores).
\item \textbf{Analysis:} 5 min (cleaning, engineering, splitting).
\end{itemize}

\artin

\artinpart{Hardware}
Parsing 1.37M Darshan logs benefits from parallel I/O. We used 128 CPU
cores with 256\,GB RAM. Minimum: 8 cores, 32\,GB RAM (longer runtime).

\artinpart{Software}
Same as $A_1$. Additionally requires \texttt{darshan-util} C library
(bundled with PyDarshan 3.5.0).

\artinpart{Datasets / Inputs}
\begin{itemize}
\item 1{,}397{,}293 \texttt{.darshan} files from ALCF Polaris
(DOI: 10.5281/zenodo.15052603).
\item Date range: 2024-04-24 to 2026-02-24 ($\sim$22 months).
\item Anonymized (uid, exe, file suffixes hashed; mount points preserved).
\end{itemize}

\artinpart{Installation and Deployment}
No additional installation beyond $A_1$. Download dataset:
\begin{verbatim}
# Download from Zenodo (requires ~50 GB)
wget https://zenodo.org/records/15052603
\end{verbatim}

\artcomp
\begin{description}
\item[$T_1$:] Batch parsing via \texttt{src/data/parse\_darshan.py}
using 7 Darshan counter aggregation rules (SUM, MAX, MIN\_NONZERO,
LAST\_VALUE, TOP-4 MERGE, CONDITIONAL, ZEROED).
99.995\% parse success rate (75 failures out of 1.4M).
\item[$T_2$:] Four-stage cleaning: require POSIX module, min duration
$\geq$\,1s, min bytes $\geq$\,1024, min ops $\geq$\,1. Reduces
1{,}397{,}216 to 131{,}151 rows.
\item[$T_3$:] Feature engineering: 39 derived features (ratios,
bandwidth, dominant access size). Normalization: 14 feature groups
with group-specific transforms.
\item[$T_4$:] Heuristic labeling via vectorized Drishti rules
(30 insight codes mapped to 8 dimensions).
\end{description}

Scripts: \texttt{scripts/reproduce\_all.sh --step 1-5}.

\artout
\begin{itemize}
\item \texttt{data/processed/production/raw\_features.parquet}:
1{,}397{,}216 rows $\times$ 156 columns (152\,MB).
\item \texttt{data/processed/production/engineered\_features.parquet}:
131{,}151 rows $\times$ 195 columns (140\,MB).
\item \texttt{data/processed/heuristic\_labels.parquet}:
131{,}151 rows $\times$ 40 columns.
\item Preprocessing figures in \texttt{paper/figures/preprocessing/}.
\end{itemize}

% ===========================================================
%  A3: Benchmark ground-truth dataset and KB
% ===========================================================
\newartifact

\artrel
Artifact $A_3$ provides the benchmark ground-truth labels ($C_2$) and
the Knowledge Base used by the LLM pipeline ($C_3$) and closed-loop
validation ($C_4$). The 623 labeled samples serve as the primary
evaluation dataset for all ML results.

\artexp
\begin{itemize}
\item 623 benchmark samples across 6 suites (IOR: 147, mdtest: 120,
DLIO: 60, custom: 89, h5bench: 123, HACC-IO: 84) with 100\%
automated verification pass rate.
\item All 8 bottleneck dimensions covered by at least 2 benchmark suites.
\item Iterative-stratified split: 187 dev / 436 test.
\item Closed-loop speedups: 3.6--39.0$\times$ on 4 IOR pairs
(geometric mean 7.2$\times$).
\end{itemize}

\arttime
\begin{itemize}
\item \textbf{Setup:} 15 min (SLURM scripts and benchmark binaries).
\item \textbf{Execution:} 180 min (726 SLURM jobs with 3 repetitions
each, sequential dependency chains).
Closed-loop validation: 30 min (8 SLURM jobs: 4 before + 4 after).
\item \textbf{Analysis:} 10 min (feature extraction, label assignment,
verification).
\end{itemize}

\artin

\artinpart{Hardware}
NCSA Delta: AMD EPYC 7763, 128 cores/node, 256\,GB DDR4, HPE
Slingshot~11 interconnect (100\,Gbps/node). Lustre parallel filesystem
(\texttt{/work}, 12 OSTs). SLURM job scheduler.

Lustre stripe configuration is critical for reproducibility:
bottleneck scenarios use \texttt{lfs setstripe -c 1} (single OST);
healthy scenarios use \texttt{lfs setstripe -c -1} (all OSTs).

\artinpart{Software}
\begin{itemize}
\item IOR 3.3.0, mdtest (bundled with IOR), DLIO (latest pip),
h5bench 1.4, HACC-IO (custom build).
\item Darshan 3.4.6 (built from source with bzip2 support).
\item Cray MPICH (default Delta module).
\item Custom Python benchmarks: mpi4py + LD\_PRELOAD for Darshan.
\end{itemize}

\artinpart{Datasets / Inputs}
Benchmark sweep configurations are specified in
\texttt{configs/benchmarks.yaml} (48 scenario configurations).
Each scenario deterministically induces a known I/O pattern;
labels are assigned by construction at job submission time,
not derived from Darshan counters (avoiding circularity).

Pre-computed benchmark features and labels are included in the
repository under \texttt{data/processed/benchmark/} for
reproduction without re-running benchmarks.

\artinpart{Installation and Deployment}
Benchmark execution requires an HPC system with Lustre and SLURM.
SLURM scripts are provided in \texttt{benchmarks/*/} directories.
For reproduction without HPC access, use the pre-computed
features in \texttt{data/processed/benchmark/}.

\artcomp
\begin{description}
\item[$T_1$:] Run benchmark sweeps (726 SLURM jobs across 6 suites).
Each job writes its configuration-based label to stdout.
Dependencies: $T_1$ is independent.
\item[$T_2$:] Extract Darshan features from benchmark logs using
the production pipeline (\texttt{src/data/parse\_darshan.py}).
Dependencies: $T_1 \rightarrow T_2$.
\item[$T_3$:] Assign labels from SLURM stdout, verify against
Darshan counter signatures via
\texttt{scripts/verify\_all\_ground\_truth.py}.
Dependencies: $T_2 \rightarrow T_3$.
\item[$T_4$:] Build Knowledge Base: for each verified benchmark
entry, store \{Darshan signature, benchmark parameters, source
code pattern, recommended fix, measured effect\} as JSON.
Dependencies: $T_3 \rightarrow T_4$.
\item[$T_5$:] Closed-loop validation: run ``before'' IOR job
(bottleneck configuration), apply LLM-recommended fix, run
``after'' IOR job, compare bandwidth.
Dependencies: $T_4 \rightarrow T_5$.
\end{description}

\artout
\begin{itemize}
\item \texttt{data/processed/benchmark/benchmark\_features.parquet}:
623 rows $\times$ 195 columns with ground-truth labels.
\item \texttt{data/knowledge\_base/*.json}: 623 KB entries.
\item \texttt{results/closed\_loop/closed\_loop\_results.json}:
Before/after bandwidth measurements for 4 IOR pairs.
\item \texttt{scripts/verify\_all\_ground\_truth.py}: Automated
verification (100\% pass rate on all 3{,}247 Darshan files).
\end{itemize}

%%%%%%%%%%%%%%%%%%%%%%%%%%%%%%%%%%%%%%%%%%%%%%%%%%%%%%%%
\section{Software Dependencies}
%%%%%%%%%%%%%%%%%%%%%%%%%%%%%%%%%%%%%%%%%%%%%%%%%%%%%%%%

All dependencies are specified in \texttt{environment.yml} (conda)
and \texttt{requirements.txt} (pip, pinned versions).

\vspace{1ex}
\begin{center}
\begin{tabular}{lll}
\toprule
Package & Version & Purpose \\
\midrule
Python      & 3.9     & Runtime \\
XGBoost     & 2.1.4   & ML classifier \\
LightGBM    & 4.5.0   & ML classifier \\
scikit-learn & 1.6.1  & ML utilities, metrics \\
SHAP        & 0.49.1  & Feature attribution \\
Cleanlab    & 2.7.1   & Label noise detection \\
pandas      & 2.3.3   & Data processing \\
NumPy       & 1.26.4  & Numerical computing \\
PyDarshan   & 3.5.0   & Darshan log parsing \\
sentence-transformers & 2.2+ & Embedding for RAG \\
faiss-cpu   & 1.7+    & Vector similarity search \\
openai      & (latest) & OpenRouter API client \\
\bottomrule
\end{tabular}
\end{center}

\vspace{1ex}
\noindent\textbf{LLM API access:} The LLM component requires an
OpenRouter API key (\texttt{OPENROUTER\_API\_KEY} environment variable)
for live inference with Claude~3.5~Sonnet, GPT-4o, and Llama~3.1~70B.
All LLM outputs are cached in \texttt{data/llm\_cache/} as JSON files,
enabling full reproduction of Tables~3--4 and Figures~7--8 without
API access.

\noindent\textbf{Target badges:} Artifacts Available (Zenodo DOI) +
Artifacts Evaluated---Functional.

\end{document}
